\section{Gated Attention: mitigación del Attention Sink}

\subsection{Idea central}

Gated Attention introduce un \textbf{mecanismo de compuerta} que permite al modelo, en ciertos Heads o posiciones, «no asignar atención», evitando así que se vea forzado a verter la atención sobre el primer token:

$$
\text{GatedAttention}(Q, K, V) = g \odot \text{softmax}\left(\frac{QK^T}{\sqrt{d_k}}\right) V
$$

donde $g$ es un vector de compuerta que puede aprenderse y que puede escalar la salida de atención de ciertas posiciones hasta acercarla a 0.

\subsection{Resumen de la sección}

Gated Attention, mediante el mecanismo de compuerta, le da al modelo «la posibilidad de no atender», y así mitiga el problema del Attention Sink que provoca la normalización forzada del Softmax.

